\documentclass{article}
\usepackage{homework}

\config{分析\ -\ 0}{2026年春季}{1}
\renewcommand{\i}{\mathrm{i}}

\begin{document}

\begin{problem}
如果$ r\neq 0 $是有理数而$ x $为无理数, 证明: $ r+x $与$ rx $均不为有理数.
\end{problem}

\begin{proof}
采用反证法. \par 
若$ p:=r+x $为有理数, 则$ x=p-r $是两个有理数的差, 也为有理数, 矛盾; \par 
若$ q:=rx $为有理数, 则$ x=q/r $是两个有理数的商, 也为有理数, 矛盾.
\end{proof}

\begin{problem}
设$ E $是某有序集的非空子集, 又设$ \alpha $为$ E $的下界, $ \beta $为$ E $的上界. 证明:$ \alpha\le\beta $.
\end{problem}

\begin{proof}
由$ E $非空, 知存在$ x\in E $. 又$ \alpha $为$ E $的下界, $ \beta $为$ E $的上界, 故$ \alpha\le x\le\beta $, 从而$ \alpha\le\beta $.
\end{proof}

\begin{problem}
设$ A $为非空实数集, 下有界. 令$ -A=\{-x:\ x\in A\} $. 求证: $ \inf A=-\sup(-A) $.
\end{problem}

\begin{proof}
由$ A $非空且下有界, 知$ \inf A $存在. 对于任意$ x\in A $, 有$ -x\in -A $, 且$ -x\le -\inf A $. 因此$ -\inf A $是$ -A $的上界, 从而$ \sup(-A)\le -\inf A $. \par
另一方面, 对于任意的$ x\in A $, 有$ -x\in -A $, 故$ -x\le \sup(-A) $, 即$ x\ge-\sup(-A) $. 因此$ -\sup(-A) $是$ A $的下界, 从而$ \inf A\ge -\sup(-A) $. \par
综上所述, $ \inf A=-\sup(-A) $.
\end{proof}

\begin{problem}
设$ z=a+b\i,\ w=c+d\i $. 若$ a<c $, 或$ a=c $且$ b<d $, 则规定$ z<w $. 证明: 这能使得所有复数的集合成为一个有序集. 这种序关系有没有最小上界性?
\end{problem}

\begin{solution}
(1) 证明“$<$”是一个序关系. \par
注意到$ z<w $的必要条件是$ a\le c $. \par
为了方便起见, 在这里给出$ z\not<w $(即$ w<z $或$ w=z $)的正面叙述: $ a>c $, 或$ a=c $且$ b\ge d $. 注意到$ z\not<w $的必要条件是$ a\ge c $. \par
\begin{enumerate}[label=(\roman*)]
\item 先证明: 若$ z,w\in\mathbb{C} $, 则$ z<w,\ z=w,\ w<z $有且仅有一个成立.\par 
\begin{enumerate}[label=\textbullet]
\item 若$ z<w $, 则$ a<c $, 此时不可能有$ z=w $或$ w<z $; 或者$ a=c $且$ b<d $, 因此不可能有$ z=w $或者$ w<z $.\par 
\item 若$ w<z $, 同理不可能有$ z<w $或$ z=w $.\par 
\item 若$ z=w $, 则$ a=c $且$ b=d $, 故不可能有$ b<d $或$ d>b $, 即不可能有$ z<w $或$ w<z $.\par 
\item 而若$ z<w,\ z=w,\ w<z $均不成立, 我们对$ a $与$ c $的大小关系分类讨论. 若$ a<c $, 这与$ z\not<w $矛盾; 若$ a>c $, 这与$ w\not<z $矛盾; 若$ a=c $, 则$ b\ge d $与$ z\not<w $矛盾, $ b\le d $与$ w\not<z $矛盾. \par
\end{enumerate}
因此, $ z<w,\ z=w,\ w<z $有且仅有一个成立.\par
\item 证明: 若$ u,v,w\in\mathbb{C} $, 且$ u<v,\ v<w $, 则$ u<w $. \par
由$ u<v $知$ a<c $, 或$ a=c $且$ b<d $; 由$ v<w $知$ c<e $, 或$ c=e $且$ d<f $ (此处记$ w=e+f\i $). 注意到一定有$ a\le c\le e $. 若$ a<e $, 则已经有$ u<w $; 若$ a=e $, 则$ a=c=e $, 因此$ b<d<f $, 从而$ u<w $.\par
\end{enumerate}\par
综上所述, “$<$”是一个序关系. \par
(2) 证明配备了这个序关系的$ \mathbb{C} $没有最小上界性. 设$ A=\{z=a+b\i:\ b<0 \} $. 则$ A $的上界为$ 0 $, 但$ A $没有最小上界, 因为对于任意的实数$ x=x+0\i $, $ x $均为$ A $的上界. 因此命题成立.
\end{solution}

\end{document}